%!TEX root = ../../main.tex
\section{측정 모형으로서의 교전 정의}
\label{sec:rel-validity}

``한타''는 플레이어의 용어이지 측정 단위가 아니다. 공식적인 정의는 없고 커뮤니티 위키는 여러 챔피언이 한 곳에서 싸우는 것으로 비형식적으로 기술한다 \cite{lolwiki2026terminology}. 존재하는 형식적 정의들도 무엇이 전투를 만드는지에 관해 서로 다르다. Schubert 등의 encounter는 상대 팀의 영웅들이 전투 범위 안에 들어올 때 시작하고 \cite{schubert2016encounter}, Ke 등의 team fight는 한쪽에 두 명 이상이 있고 킬이 적어도 하나 있는 encounter이며 \cite{ke2022}, OpenDota의 라벨은 죽음만을 요구한다 \cite{opendota2018processteamfights,tot2021camera}.

따라서 본 논문은 교전을 측정 모형(measurement modelling)의 의미에서 \emph{조작화된 구성개념}으로 다룬다 \cite{jacobs2021measurement}. 관측 가능한 텔레메트리를 관측 불가능한 개념의 사례로 사상하는 규칙이되, 그 가정을 명시하고 신뢰도와 타당도를 주장이 아니라 검사로 다룬다. Cronbach와 Meehl \cite{cronbach1955construct}의 구성 타당성 개념에 따라 내용·수렴·판별 타당성과 검사--재검사 신뢰도, 예측 타당성을 제 \ref{ch:engagement} 장 끝에서 항목별로 요약한다.
